\documentclass{tietreport}

% ============================================================
% PACKAGES
% ============================================================

\usepackage{graphicx}
\usepackage{float}
\usepackage{array}
\usepackage{longtable}
\usepackage{booktabs}
\usepackage{tabularx}
\usepackage{enumitem}
\usepackage{csquotes}

% ============================================================
% BIBLIOGRAPHY
% ============================================================

\usepackage[
    backend=biber,
    style=alphabetic,
    sorting=nyt
]{biblatex}

\addbibresource{references.bib}

% ============================================================
% GENERAL FORMATTING
% ============================================================

% Better table columns
\newcolumntype{Y}{>{\raggedright\arraybackslash}X}

% Longtable alignment
\setlength{\LTleft}{0pt}
\setlength{\LTright}{0pt}

% Paragraph formatting
\setlength{\parindent}{0pt}
\setlength{\parskip}{6pt}

% Compact but readable lists
\setlist[itemize]{itemsep=2pt, topsep=4pt}
\setlist[enumerate]{itemsep=2pt, topsep=4pt}

% ============================================================
% DOCUMENT
% ============================================================

\begin{document}

% ============================================================
% TITLE PAGE
% ============================================================

\begin{titlepage}

\thispagestyle{empty}

\begin{center}


{\large\bfseries
UCS 503\\
Software Engineering Practice Project Report
}

\vspace{2.0cm}

{\Huge\bfseries CampusCare}

\vspace{0.6cm}

{\large
An Integrated Digital Healthcare Platform\\
for Thapar's Health Center
}

\vspace{2.0cm}

{\large\bfseries Submitted by}

\vspace{0.6cm}

\begin{tabular}{ll}
1024030107 Lavdeep Singh \\
1024030109 Harshveer Singh \\
1024030110 Damanbir Singh \\
1024030111 Kriti
\end{tabular}

\vspace{1.2cm}

{\large
BE Third Year\\
Computer Science and Engineering
}

\vspace{1.2cm}

{\large\bfseries Under the Guidance of}

\vspace{0.5cm}

{\large
Ms. Anushka
}

\vfill

{\large
Department of Computer Science and Engineering\\
Thapar Institute of Engineering and Technology\\
Patiala
}

\end{center}

\end{titlepage}

% ============================================================
% TABLE OF CONTENTS
% ============================================================

\tableofcontents

\clearpage


% ============================================================
% CHAPTER 1
% INTRODUCTION
% ============================================================

\chapter{Introduction}

\section{Background}

CampusCare is an integrated digital healthcare platform designed for
campus health centers. The system brings the major activities of a
routine student healthcare visit together within a single platform.

In a conventional campus health-center environment, student
registration, doctor queues, consultation records and prescriptions
may be handled through separate or partially manual processes. This
can make it difficult to maintain consistent records and can increase
the amount of administrative work required from healthcare-center
staff.

CampusCare addresses this problem through a centralized,
role-based system. The platform connects student identification and
check-in with doctor availability, queue management, consultation,
medical history, prescription generation and medicine dispensing.

The system is divided into six major functional modules:

\begin{enumerate}
    \item Authentication and User Management
    \item Student Profile and Check-in
    \item Doctor Availability and Queue Management
    \item Consultation and Medical History
    \item Prescription and Dispensing
    \item Administration and Audit
\end{enumerate}

These modules represent the main workflow of the system and are
reflected in the Data Flow Diagrams and UML models presented later in
this report.

\section{Problem Statement}

Campus health centers need an organized method for managing student
visits, doctor availability, waiting queues, medical records and
prescriptions.

When these activities are handled manually or through disconnected
systems, several problems can occur:

\begin{itemize}
    \item student information may have to be entered or verified
    repeatedly;
    \item managing separate doctor queues can become difficult;
    \item previous medical information may take time to retrieve;
    \item consultation and prescription records may become
    disconnected;
    \item dispensing information may not be clearly linked to the
    original prescription;
    \item administrative activities may be difficult to track.
\end{itemize}

CampusCare is intended to provide a centralized solution that manages
the routine healthcare workflow from student check-in through
consultation, prescription and dispensing.

\section{Objectives}

The main objectives of CampusCare are:

\begin{enumerate}
    \item To provide secure authentication and role-based access.
    
    \item To maintain student profile and identification information.
    
    \item To provide an organized student check-in process.
    
    \item To manage doctor availability and doctor-specific queues.
    
    \item To provide authorized access to student medical history.
    
    \item To allow doctors to record consultation information.
    
    \item To allow doctors to create structured prescriptions.
    
    \item To maintain separate prescription and dispensing records.
    
    \item To provide administrative controls and audit information.
    
    \item To provide a foundation for RFID/NFC-based identification
    where suitable hardware and browser support are available.
\end{enumerate}

\section{Scope}

CampusCare focuses on the routine healthcare-management activities
performed within a campus health center.

The core scope includes:

\begin{itemize}
    \item authentication and user management;
    \item student profiles;
    \item student identification and check-in;
    \item doctor availability;
    \item doctor-specific queue management;
    \item consultation management;
    \item medical-history management;
    \item prescription generation;
    \item medicine dispensing;
    \item administration;
    \item audit information.
\end{itemize}

RFID/NFC-based identification, advanced real-time functionality,
medication reminders and advanced reporting can be considered
extensions to the core system.

Inventory management is outside the current project scope.

\section{Users of the System}

CampusCare provides different interfaces and permissions for the
following roles.

\subsection{Student}

Students use the system for authentication, identification,
check-in and accessing information relevant to their visits.

\subsection{Doctor}

Doctors manage their availability, access authorized student medical
history, conduct consultations and create prescriptions.

\subsection{Receptionist}

Receptionists assist students with identification and check-in and
manage routine queue-related operations.

\subsection{Pharmacist}

Pharmacists review valid prescriptions and record medicine dispensing.

\subsection{Administrator}

Administrators manage user accounts, system configuration, doctor
availability, identification credentials and audit information.


% ============================================================
% CHAPTER 2
% REQUIREMENTS ANALYSIS
% ============================================================

\chapter{Requirements Analysis}

\section{Functional Requirements}

The functional requirements of CampusCare are listed below.

\begin{longtable}{p{1.3cm}p{4.0cm}p{7.0cm}}
\toprule
\textbf{ID} & \textbf{Requirement} & \textbf{Description} \\
\midrule

FR-01 &
Authentication &
The system shall authenticate users and provide access according to
their assigned role. \\

FR-02 &
Student Profile &
The system shall maintain student profile and identification
information. \\

FR-03 &
Student Check-in &
The receptionist shall be able to verify and check in a student. \\

FR-04 &
Doctor Availability &
The system shall maintain doctor availability information. \\

FR-05 &
Queue Management &
The system shall maintain separate queues for different doctors. \\

FR-06 &
Queue Status &
The system shall maintain queue position and status information. \\

FR-07 &
Medical History &
An authorized doctor shall be able to access relevant student medical
history. \\

FR-08 &
Consultation &
The doctor shall be able to record the reason for consultation,
diagnosis and clinical notes. \\

FR-09 &
Prescription &
The doctor shall be able to create a structured prescription linked
to the consultation. \\

FR-10 &
Dispensing &
The pharmacist shall be able to record medicine dispensing against a
valid prescription. \\

FR-11 &
Administration &
The administrator shall be able to manage users, configuration and
operational information. \\

FR-12 &
RFID/NFC &
The system may support RFID/NFC-based identification where suitable
hardware and browser support are available. \\

FR-13 &
Audit &
Important administrative and system actions shall be recorded for
auditing. \\

FR-14 &
No-show Handling &
The queue system shall support handling of students who do not
respond when called. \\

FR-15 &
Queue Redirection &
Authorized staff shall be able to redirect a student to an
appropriate doctor queue when required. \\

\bottomrule
\end{longtable}

\section{Non-Functional Requirements}

\subsection{Security}

CampusCare shall restrict system functionality and healthcare
information according to user roles. Authentication and
authorization are therefore important parts of the system.

\subsection{Reliability}

Student, consultation, prescription and dispensing information should
remain consistent and available during normal system operation.

\subsection{Usability}

The system should provide simple and understandable interfaces for
students and healthcare-center staff.

\subsection{Maintainability}

The application should be modular so that individual components can
be modified and maintained without unnecessary changes to unrelated
components.

\subsection{Performance}

Routine operations such as login, check-in, queue operations and
record retrieval should provide reasonable response times under
normal campus health-center usage.

\subsection{Auditability}

Important administrative and system activities should be traceable
through audit information.

\section{Technology Stack}

The CampusCare technology stack consists of:

\begin{itemize}
    \item \textbf{Frontend:} React with Vite
    \item \textbf{Styling:} Tailwind CSS
    \item \textbf{Backend:} FastAPI
    \item \textbf{API:} REST API
    \item \textbf{Database:} MySQL
    \item \textbf{Database Management:} SQL with
    version-controlled migrations
    \item \textbf{Identification:} Staff-assisted check-in with
    planned RFID/NFC support
    \item \textbf{Testing:} Unit and integration testing
\end{itemize}


% ============================================================
% CHAPTER 3
% SYSTEM DESIGN AND ARCHITECTURE
% ============================================================

\chapter{System Design and Architecture}

\section{System Architecture}

CampusCare follows a client-server architecture consisting of a
frontend, backend and relational database.

The frontend is responsible for providing role-specific interfaces
for students, doctors, receptionists, pharmacists and administrators.

The FastAPI backend handles authentication, authorization, business
logic, queue operations, consultation processing, prescriptions and
administrative operations.

MySQL provides persistent storage for the system data.

The overall architecture can be represented as:

\begin{center}
\textbf{React Frontend}
$\rightarrow$
\textbf{FastAPI REST API}
$\rightarrow$
\textbf{Application Logic}
$\rightarrow$
\textbf{MySQL Database}
\end{center}

This separation allows the user interface, application logic and
database layer to be developed and maintained independently.

% ------------------------------------------------------------
% CONTEXT DIAGRAM
% ------------------------------------------------------------

\section{Context Diagram}

The Context Diagram represents CampusCare as a single system and
shows its interaction with the external actors.

\begin{figure}[H]
    \centering
    \includegraphics[
        width=0.90\textwidth,
        keepaspectratio
    ]{levle000.jpeg}
    \caption{CampusCare Context Diagram}
    \label{fig:context}
\end{figure}

The main external actors are:

\begin{itemize}
    \item Student
    \item Doctor
    \item Receptionist
    \item Pharmacist
    \item Administrator
\end{itemize}

The diagram provides a high-level view of the information exchanged
between these actors and the CampusCare system.

% ------------------------------------------------------------
% LEVEL 1 DFD
% ------------------------------------------------------------

\section{Level-1 Data Flow Diagram}

The Level-1 DFD decomposes CampusCare into six major processes.

\begin{figure}[H]
    \centering
    \includegraphics[
        width=0.95\textwidth,
        keepaspectratio
    ]{level1 se (1).png}
    \caption{CampusCare Level-1 Data Flow Diagram}
    \label{fig:level1}
\end{figure}

The six Level-1 processes are:

\begin{enumerate}
    \item \textbf{1.0 Authentication and User Management}
    \item \textbf{2.0 Student Profile and Check-in}
    \item \textbf{3.0 Doctor Availability and Queue Management}
    \item \textbf{4.0 Consultation and Medical History}
    \item \textbf{5.0 Prescription and Dispensing}
    \item \textbf{6.0 Administration and Audit}
\end{enumerate}

The Level-2 diagrams below provide a more detailed decomposition of
each of these processes.

% ------------------------------------------------------------
% LEVEL 2
% ------------------------------------------------------------

\section{Level-2 Data Flow Diagrams}

\subsection{1.0 Authentication and User Management}

The first Level-2 DFD decomposes authentication and user management
into smaller activities such as authentication, authorization and
user-account management.

\begin{figure}[H]
    \centering
    \includegraphics[
        width=0.95\textwidth,
        keepaspectratio
    ]{level2_1.png}
    \caption{Level-2 DFD for Authentication and User Management}
    \label{fig:l2-authentication}
\end{figure}

\subsection{2.0 Student Profile and Check-in}

This Level-2 DFD describes student identification, profile
verification and the check-in process.

\begin{figure}[H]
    \centering
    \includegraphics[
        width=0.95\textwidth,
        keepaspectratio
    ]{level2_2.png}
    \caption{Level-2 DFD for Student Profile and Check-in}
    \label{fig:l2-checkin}
\end{figure}

\subsection{3.0 Doctor Availability and Queue Management}

This Level-2 DFD describes doctor availability, queue creation,
queue position, queue status, no-show handling and queue redirection.

\begin{figure}[H]
    \centering
    \includegraphics[
        width=0.95\textwidth,
        keepaspectratio
    ]{level2_3.png}
    \caption{Level-2 DFD for Doctor Availability and Queue Management}
    \label{fig:l2-queue}
\end{figure}

\subsection{4.0 Consultation and Medical History}

This Level-2 DFD describes authorized access to student medical
history and the recording of consultation information.

\begin{figure}[H]
    \centering
    \includegraphics[
        width=0.95\textwidth,
        keepaspectratio
    ]{level2_4.png}
    \caption{Level-2 DFD for Consultation and Medical History}
    \label{fig:l2-consultation}
\end{figure}

\subsection{5.0 Prescription and Dispensing}

This Level-2 DFD describes prescription creation by the doctor and
medicine dispensing by the pharmacist.

The prescription and dispensing records are kept separate. The
prescription represents the doctor's original instruction, while the
dispensing record represents the pharmacist's action.

\begin{figure}[H]
    \centering
    \includegraphics[
        width=0.95\textwidth,
        keepaspectratio
    ]{level2_5.png}
    \caption{Level-2 DFD for Prescription and Dispensing}
    \label{fig:l2-prescription}
\end{figure}

\subsection{6.0 Administration and Audit}

This Level-2 DFD describes administrative operations such as user
management, doctor availability management, identification
credential management, record retrieval and audit operations.

\begin{figure}[H]
    \centering
    \includegraphics[
        width=0.95\textwidth,
        keepaspectratio
    ]{Level6_diag.png}
    \caption{Level-2 DFD for Administration and Audit}
    \label{fig:l2-admin}
\end{figure}

% ------------------------------------------------------------
% UML
% ------------------------------------------------------------

\section{UML Design}

UML diagrams provide complementary views of the CampusCare system.
Four UML diagrams are included:\newline

1. Use Case Diagram \newline
2. Activity Diagram \newline
3. Class Diagram\newline
4. Sequence Diagram\newline
These diagrams describe the system from the functional, behavioral,
structural and interaction perspectives.

% ------------------------------------------------------------
% USE CASE
% ------------------------------------------------------------

\subsection{Use Case Diagram}

The Use Case Diagram represents the interactions between the
different CampusCare actors and the system.

The main actors are Student, Doctor, Receptionist, Pharmacist and
Administrator.

Major use cases include authentication, student check-in, queue
management, consultation, medical-history access, prescription
creation, medicine dispensing and administrative management.

\begin{figure}[H]
    \centering
    \includegraphics[
        width=0.95\textwidth,
        keepaspectratio
    ]{User_case_dig.png}
    \caption{CampusCare Use Case Diagram}
    \label{fig:usecase}
\end{figure}

% ------------------------------------------------------------
% ACTIVITY
% ------------------------------------------------------------

\subsection{Activity Diagram}

The Activity Diagram represents the normal workflow of a student
health-center visit.

The process begins with authentication and student identification,
followed by check-in, queue assignment, consultation and, when
required, prescription and dispensing.

\begin{figure}[H]
    \centering
    \includegraphics[
        width=0.92\textwidth,
        keepaspectratio
    ]{Activity_diag.jpeg}
    \caption{CampusCare Activity Diagram}
    \label{fig:activity}
\end{figure}

The activity model focuses on the normal healthcare workflow and does
not include any emergency or ambulance workflow.

% ------------------------------------------------------------
% CLASS
% ------------------------------------------------------------

\subsection{Class Diagram}

The Class Diagram represents the main structural entities of
CampusCare and their relationships.

The model covers users and roles, student information, doctor
availability, visits, queue entries, consultations, prescriptions,
prescription items, dispensing records, RFID/NFC credentials and
audit information.

\begin{figure}[H]
    \centering
    \includegraphics[
        width=0.95\textwidth,
        keepaspectratio
    ]{class_daig.jpeg}
    \caption{CampusCare Class Diagram}
    \label{fig:class}
\end{figure}

The class model is focused on the current CampusCare database and
does not contain emergency or ambulance classes.

% ------------------------------------------------------------
% SEQUENCE
% ------------------------------------------------------------

\subsection{Sequence Diagram}

The Sequence Diagram represents the interaction between the user,
CampusCare application services and the MySQL database during a
typical healthcare visit.

The sequence covers authentication, student check-in, queue
processing, consultation, prescription creation and dispensing.

\begin{figure}[H]
    \centering
    \includegraphics[
        width=0.95\textwidth,
        keepaspectratio
    ]{Sequence_diag.jpeg}
    \caption{CampusCare Sequence Diagram}
    \label{fig:sequence}
\end{figure}


% ============================================================
% DATABASE DESIGN
% ============================================================

\section{Database Design}

CampusCare uses MySQL as its relational database.

The database is organized around the lifecycle of a routine student
visit:

\begin{center}
\textbf{Student}
$\rightarrow$
\textbf{Visit}
$\rightarrow$
\textbf{Queue Entry}
$\rightarrow$
\textbf{Consultation}
$\rightarrow$
\textbf{Prescription}
$\rightarrow$
\textbf{Dispensing}
\end{center}

The major conceptual entities include:

\begin{itemize}
    \item User
    \item Student
    \item Doctor
    \item Receptionist
    \item Pharmacist
    \item Administrator
    \item RFID Card
    \item Doctor Availability
    \item Visit
    \item Queue Entry
    \item Consultation
    \item Prescription
    \item Prescription Item
    \item Dispensing Record
    \item Audit Log
\end{itemize}

The exact implementation of these entities is maintained through the
MySQL database schema and version-controlled migration files.

\section{Prescription and Dispensing Separation}

A major database-design decision is to keep prescriptions and
dispensing as separate records.

A prescription records what the doctor has prescribed. A dispensing
record records what the pharmacist has actually dispensed.

This separation prevents the dispensing process from silently
changing the doctor's original prescription and provides a clearer
record of the healthcare workflow.

\section{Database Migration Strategy}

Database schema changes are managed using version-controlled
migration files.

This allows changes to the MySQL schema to be tracked and reproduced
across development environments.


% ============================================================
% CHAPTER 4
% SYSTEM IMPLEMENTATION
% ============================================================

\chapter{System Implementation}

\section{Frontend Implementation}

The CampusCare frontend is developed using React with Vite.
Tailwind CSS is used for styling and interface development.

The frontend provides role-specific screens for the different
CampusCare users.

\subsection{Login Page}

The login page provides the entry point for user authentication.

\begin{figure}[H]
    \centering
    \includegraphics[
        width=0.88\textwidth,
        keepaspectratio
    ]{login_page.png}
    \caption{CampusCare Login Page}
    \label{fig:login}
\end{figure}

\subsection{Landing Page}

The landing page provides access to functionality available to the
authenticated user.

\begin{figure}[H]
    \centering
    \includegraphics[
        width=0.88\textwidth,
        keepaspectratio
    ]{Lnding_page.jpeg}
    \caption{CampusCare Landing Page}
    \label{fig:landing}
\end{figure}

\subsection{Queue Management Console}

The queue management console allows staff to monitor and manage
doctor-specific queues.

\begin{figure}[H]
    \centering
    \includegraphics[
        width=0.90\textwidth,
        keepaspectratio
    ]{live_queue_console.png}
    \caption{CampusCare Queue Management Console}
    \label{fig:queue-console}
\end{figure}

\subsection{Doctor Dashboard -- Patient Record}

The doctor dashboard provides access to authorized student records
and relevant medical information.

\begin{figure}[H]
    \centering
    \includegraphics[
        width=0.90\textwidth,
        keepaspectratio
    ]{doctor_dashboard_patient_record.jpeg}
    \caption{Doctor Dashboard -- Patient Record}
    \label{fig:doctor-record}
\end{figure}

\subsection{Doctor Dashboard -- Prescription}

The doctor dashboard also provides functionality for creating
structured prescriptions.

\begin{figure}[H]
    \centering
    \includegraphics[
        width=0.90\textwidth,
        keepaspectratio
    ]{doctor_dashboard_prescription.jpeg}
    \caption{Doctor Dashboard -- Prescription}
    \label{fig:doctor-prescription}
\end{figure}

\section{Backend Implementation}

The backend is implemented using FastAPI and exposes REST APIs for
communication between the frontend and application services.

The backend is responsible for:

\begin{itemize}
    \item authentication and authorization;
    \item user and role management;
    \item student profile management;
    \item student check-in;
    \item visit creation;
    \item doctor availability;
    \item queue management;
    \item medical-history access;
    \item consultation recording;
    \item prescription management;
    \item dispensing records;
    \item administrative operations;
    \item audit information.
\end{itemize}

\section{Role-Based Access Control}

Role-based access control separates the capabilities of students,
doctors, receptionists, pharmacists and administrators.

This ensures that users are given access only to the functions
required for their responsibilities.


% ============================================================
% CHAPTER 5
% PROJECT DEVELOPMENT AND MILESTONES
% ============================================================

\chapter{Project Development and Milestones}

\section{Development Phases}

The development of CampusCare can be organized into the following
phases.

\begin{longtable}{p{1.2cm}p{4.0cm}p{7.0cm}}
\toprule
\textbf{Phase} & \textbf{Area} & \textbf{Major Work} \\
\midrule

1 &
Authentication and Profiles &
Authentication, user roles and student profile management. \\

2 &
Check-in and Queue &
Student check-in, doctor availability and queue management. \\

3 &
Consultation and Medical History &
Medical-history access and consultation recording. \\

4 &
Prescription and Dispensing &
Prescription creation and separate dispensing records. \\

5 &
Administration and Audit &
Administrative controls, configuration and audit information. \\

6 &
Extensions &
RFID/NFC support, advanced real-time functionality, reminders,
analytics and additional reporting. \\

\bottomrule
\end{longtable}

\section{Core System Workflow}

The normal CampusCare workflow is:

\begin{enumerate}

    \item The user authenticates into the system.

    \item The student's identity is verified.

    \item The receptionist checks in the student.

    \item A visit is created.

    \item The student is added to the appropriate doctor's queue.

    \item The doctor calls the next student.

    \item The doctor accesses the student's authorized medical history.

    \item The doctor conducts and records the consultation.

    \item The doctor creates a prescription when required.

    \item The pharmacist reviews the prescription.

    \item The pharmacist records medicine dispensing.

    \item The visit information remains available for future
    authorized access.

\end{enumerate}


% ============================================================
% CHAPTER 6
% TESTING AND VALIDATION
% ============================================================

\chapter{Testing and Validation}

\section{Testing Strategy}

CampusCare is tested through functional, integration and
role-based testing.

The main areas considered during testing are:

\begin{itemize}
    \item authentication;
    \item student check-in;
    \item doctor availability;
    \item queue management;
    \item medical-history access;
    \item consultation recording;
    \item prescription creation;
    \item medicine dispensing;
    \item role authorization;
    \item administrative auditing.
\end{itemize}

\section{Functional Test Cases}

\begin{longtable}{p{1.2cm}p{4.0cm}p{5.5cm}p{1.8cm}}
\toprule
\textbf{ID} &
\textbf{Test Case} &
\textbf{Expected Result} &
\textbf{Status} \\
\midrule

TC-01 &
Valid Login &
Valid credentials authenticate the user and provide the correct
role-based interface. &
To be recorded \\

TC-02 &
Invalid Login &
Invalid credentials are rejected. &
To be recorded \\

TC-03 &
Student Check-in &
A valid student can be checked in by the receptionist. &
To be recorded \\

TC-04 &
Queue Assignment &
The student is added to the correct doctor's queue. &
To be recorded \\

TC-05 &
Queue Status &
The current queue position and status are displayed correctly. &
To be recorded \\

TC-06 &
Doctor Availability &
Doctor availability changes are reflected in queue operations. &
To be recorded \\

TC-07 &
Medical History &
An authorized doctor can access the relevant student medical
history. &
To be recorded \\

TC-08 &
Consultation &
The doctor can record consultation information. &
To be recorded \\

TC-09 &
Prescription &
The doctor can create a structured prescription linked to the
consultation. &
To be recorded \\

TC-10 &
Dispensing &
The pharmacist can record dispensing against a valid prescription. &
To be recorded \\

TC-11 &
Role Authorization &
Users cannot access functionality outside their assigned role. &
To be recorded \\

TC-12 &
Audit &
Important administrative actions are recorded correctly. &
To be recorded \\

\bottomrule
\end{longtable}

\section{Validation}

The test cases provide a structured way to validate the major
functional requirements of CampusCare.

The status column is intentionally kept as ``To be recorded'' unless
actual test results have been measured. This avoids presenting
unverified testing results as completed results.


% ============================================================
% CHAPTER 7
% RISKS, CHALLENGES AND MITIGATIONS
% ============================================================

\chapter{Risks, Challenges and Mitigations}

\section{Security and Privacy}

CampusCare handles student healthcare information, making controlled
access important. Role-based authorization is used to limit access
according to user responsibilities.

\section{Queue Consistency}

Queue information can become inconsistent if multiple users attempt
to update the queue simultaneously.

The backend therefore acts as the authoritative source for queue
state and operations.

\section{RFID/NFC Dependency}

RFID/NFC functionality depends on suitable hardware and browser or
device support.

Staff-assisted identification remains available when RFID/NFC
hardware is not available.

\section{Data Integrity}

Prescription and dispensing information are maintained separately.

This preserves the doctor's original prescription while allowing the
pharmacist to record the dispensing action.

\section{System Availability}

CampusCare may be required during health-center operating hours.
Therefore, backend and database availability should be considered
during deployment.


% ============================================================
% CHAPTER 8
% PROJECT SCOPE AND DELIVERABLES
% ============================================================

\chapter{Project Scope and Deliverables}

\section{Core Deliverables}

The core CampusCare system includes:

\begin{itemize}

    \item Authentication and role-based access.
    
    \item Student profile management.
    
    \item Staff-assisted student identification and check-in.
    
    \item Doctor availability management.
    
    \item Doctor-specific queue management.
    
    \item Queue position and status management.
    
    \item Medical-history access for authorized doctors.
    
    \item Consultation recording.
    
    \item Structured prescription creation.
    
    \item Separate dispensing records.
    
    \item Administrative management.
    
    \item Audit information.

\end{itemize}

\section{Extended Functionality}

The following features can extend the core system:

\begin{itemize}
    \item RFID/NFC-based student identification;
    \item advanced real-time queue updates;
    \item additional queue-management functionality;
    \item additional operational reports.
\end{itemize}

\section{Stretch Functionality}

Possible future enhancements include:

\begin{itemize}
    \item medication reminders;
    \item aggregated and de-identified health trends;
    \item advanced reporting and analytics;
    \item integration with additional campus healthcare services.
\end{itemize}

\section{Out of Scope}

The following functionality is outside the current project scope:

\begin{itemize}
    \item inventory management;
    \item advanced healthcare analytics;
    \item external healthcare-system integration.
\end{itemize}


% ============================================================
% CHAPTER 9
% DISCUSSION AND LIMITATIONS
% ============================================================

\chapter{Discussion and Limitations}

\section{Discussion}

CampusCare connects the major routine healthcare-center activities
within a single digital platform.

The system links student check-in and doctor queues with consultation
records, prescriptions and dispensing. This creates a continuous
workflow instead of treating each activity as an isolated process.

A key design decision is the separation between prescriptions and
dispensing records. The prescription represents the doctor's
instruction, while the dispensing record represents the pharmacist's
action.

The DFD models describe how information moves through the system,
while the UML diagrams provide different views of functionality,
workflow, structure and interaction.

\section{Limitations}

The current project has several limitations:

\begin{itemize}

    \item RFID/NFC functionality depends on available hardware and
    browser capabilities.

    \item Advanced real-time synchronization may require additional
    infrastructure.

    \item Quantitative performance measurements require controlled
    testing under realistic workloads.

    \item Advanced analytics are outside the current core scope.

    \item Inventory management is not included.

\end{itemize}


% ============================================================
% CHAPTER 10
% CONCLUSION AND FUTURE WORK
% ============================================================

\chapter{Conclusion and Future Work}

\section{Conclusion}

CampusCare provides a centralized approach to managing routine
healthcare activities within a campus health center.

The system brings together authentication, student profiles,
check-in, doctor availability, queue management, consultation,
medical history, prescriptions, dispensing, administration and
auditing.

Role-based access separates the responsibilities of students,
doctors, receptionists, pharmacists and administrators.

The separation of prescriptions and dispensing records helps
preserve the integrity of doctor-created medical instructions while
maintaining a clear dispensing history.

The system is represented using Context and Data Flow Diagrams as well
as UML Use Case, Activity, Class and Sequence Diagrams. Together,
these models connect the requirements of CampusCare with its system
design and implementation.

\section{Future Work}

Future development can include:

\begin{itemize}
    \item broader RFID/NFC integration;
    \item improved real-time queue updates;
    \item medication reminders;
    \item aggregated and de-identified health analytics;
    \item advanced reporting;
    \item improved deployment and monitoring;
    \item integration with additional campus healthcare services.
\end{itemize}


% ============================================================
% REFERENCES
% ============================================================

\nocite{*}

\printbibliography[
    heading=bibintoc,
    title={References}
]

% ============================================================
% END OF DOCUMENT
% ============================================================

\end{document}